\begin{thebibliography}{99}
\bibitem{bib-bib1} 
\par\noindent    Amal Ahmed, Derek Dreyer, and Andreas Rossberg.  State-dependent representation independence.  In Zhong Shao and Benjamin C. Pierce, editors, Proceedings of the 36th ACM SIGPLAN-SIGACT Symposium on Principles of Programming Languages, POPL 2009, Savannah, GA, USA, January 21-23, 2009, pages 340–353. ACM, 2009.  doi:10.1145/1480881.1480925.  

\bibitem{bib-bib2} 
\par\noindent    Benedikt Ahrens, Krzysztof Kapulkin, and Michael Shulman.  Univalent categories and the Rezk completion.  Math. Struct. Comput. Sci., 25(5):1010–1039, 2015.  doi:10.1017/S0960129514000486.  

\bibitem{bib-bib3} 
\par\noindent    Pierre America and Jan J. M. M. Rutten.  Solving Reflexive Domain Equations in a Category of Complete Metric Spaces.  J. Comput. Syst. Sci., 39(3):343–375, 1989.  doi:10.1016/0022-0000(89)90027-5.  

\bibitem{bib-bib4} 
\par\noindent    Andrew W. Appel and David A. McAllester.  An indexed model of recursive types for foundational proof-carrying code.  ACM Trans. Program. Lang. Syst., 23(5):657–683, 2001.  doi:10.1145/504709.504712.  

\bibitem{bib-bib5} 
\par\noindent    Steve Awodey.  Category Theory.  Oxford Logic Guides. Oxford University Press, London, England, 2 edition, June 2010.  

\bibitem{bib-bib6} 
\par\noindent    Nick Benton, Andrew Kennedy, and Carsten Varming.  Some Domain Theory and Denotational Semantics in Coq.  In Stefan Berghofer, Tobias Nipkow, Christian Urban, and Makarius Wenzel, editors, Theorem Proving in Higher Order Logics, 22nd International Conference, TPHOLs 2009, Munich, Germany, August 17-20, 2009. Proceedings, volume 5674 of Lecture Notes in Computer Science, pages 115–130. Springer, 2009.  doi:10.1007/978-3-642-03359-9\_10.  

\bibitem{bib-bib7} 
\par\noindent    Lars Birkedal, Ales Bizjak, and Jan Schwinghammer.  Step-Indexed Relational Reasoning for Countable Nondeterminism.  Log. Methods Comput. Sci., 9(4), 2013.  doi:10.2168/LMCS-9(4:4)2013.  

\bibitem{bib-bib8} 
\par\noindent    Lars Birkedal and Aleš Bizjak.  Lecture Notes on Iris: Higher-Order Concurrent Separation Logic, 2017.  URL: http://iris-project.org/tutorial-pdfs/iris-lecture-notes.pdf.  

\bibitem{bib-bib9} 
\par\noindent    Lars Birkedal, Rasmus Ejlers Møgelberg, Jan Schwinghammer, and Kristian Støvring.  First steps in synthetic guarded domain theory: step-indexing in the topos of trees.  Log. Methods Comput. Sci., 8(4), 2012.  doi:10.2168/LMCS-8(4:1)2012.  

\bibitem{bib-bib10} 
\par\noindent    Lars Birkedal, Bernhard Reus, Jan Schwinghammer, Kristian Støvring, Jacob Thamsborg, and Hongseok Yang.  Step-indexed kripke models over recursive worlds.  In Thomas Ball and Mooly Sagiv, editors, Proceedings of the 38th ACM SIGPLAN-SIGACT Symposium on Principles of Programming Languages, POPL 2011, Austin, TX, USA, January 26-28, 2011, pages 119–132. ACM, 2011.  doi:10.1145/1926385.1926401.  

\bibitem{bib-bib11} 
\par\noindent    Lars Birkedal, Kristian Støvring, and Jacob Thamsborg.  The category-theoretic solution of recursive metric-space equations.  Theor. Comput. Sci., 411(47):4102–4122, 2010.  doi:10.1016/J.TCS.2010.07.010.  

\bibitem{bib-bib12} 
\par\noindent    Errett Bishop.  Foundations of Constructive Analysis.  Mcgraw-Hill, New York, NY, USA, 1967.  

\bibitem{bib-bib13} 
\par\noindent    Ales Bizjak, Lars Birkedal, and Marino Miculan.  A Model of Countable Nondeterminism in Guarded Type Theory.  In Gilles Dowek, editor, Rewriting and Typed Lambda Calculi - Joint International Conference, RTA-TLCA 2014, Held as Part of the Vienna Summer of Logic, VSL 2014, Vienna, Austria, July 14-17, 2014. Proceedings, volume 8560 of Lecture Notes in Computer Science, pages 108–123. Springer, 2014.  doi:10.1007/978-3-319-08918-8\_8.  

\bibitem{bib-bib14} 
\par\noindent    Ranald Clouston, Ales Bizjak, Hans Bugge Grathwohl, and Lars Birkedal.  Programming and Reasoning with Guarded Recursion for Coinductive Types.  In Andrew M. Pitts, editor, Foundations of Software Science and Computation Structures - 18th International Conference, FoSSaCS 2015, Held as Part of the European Joint Conferences on Theory and Practice of Software, ETAPS 2015, London, UK, April 11-18, 2015. Proceedings, volume 9034 of Lecture Notes in Computer Science, pages 407–421. Springer, 2015.  doi:10.1007/978-3-662-46678-0\_26.  

\bibitem{bib-bib15} 
\par\noindent    Robert Dockins.  Formalized, Effective Domain Theory in Coq.  In Gerwin Klein and Ruben Gamboa, editors, Interactive Theorem Proving - 5th International Conference, ITP 2014, Held as Part of the Vienna Summer of Logic, VSL 2014, Vienna, Austria, July 14-17, 2014. Proceedings, volume 8558 of Lecture Notes in Computer Science, pages 209–225. Springer, 2014.  doi:10.1007/978-3-319-08970-6\_14.  

\bibitem{bib-bib16} 
\par\noindent    Peter Freyd.  Algebraically complete categories.  In Aurelio Carboni, Maria Cristina Pedicchio, and Guiseppe Rosolini, editors, Category Theory, pages 95–104, Berlin, Heidelberg, 1991. Springer Berlin Heidelberg.  

\bibitem{bib-bib17} 
\par\noindent    Dan Frumin, Amin Timany, and Lars Birkedal.  Modular Denotational Semantics for Effects with Guarded Interaction Trees.  Proc. ACM Program. Lang., 8(POPL):332–361, 2024.  doi:10.1145/3632854.  

\bibitem{bib-bib18} 
\par\noindent    Pietro Di Gianantonio and Marino Miculan.  A unifying approach to recursive and co-recursive definitions.  In Herman Geuvers and Freek Wiedijk, editors, Types for Proofs and Programs, Second International Workshop, TYPES 2002, Berg en Dal, The Netherlands, April 24-28, 2002, Selected Papers, volume 2646 of Lecture Notes in Computer Science, pages 148–161. Springer, 2002.  doi:10.1007/3-540-39185-1\_9.  

\bibitem{bib-bib19} 
\par\noindent    Pietro Di Gianantonio and Marino Miculan.  Unifying recursive and co-recursive definitions in sheaf categories.  In Igor Walukiewicz, editor, Foundations of Software Science and Computation Structures, 7th International Conference, FOSSACS 2004, Held as Part of the Joint European Conferences on Theory and Practice of Software, ETAPS 2004, Barcelona, Spain, March 29 - April 2, 2004, Proceedings, volume 2987 of Lecture Notes in Computer Science, pages 136–150. Springer, 2004.  doi:10.1007/978-3-540-24727-2\_11.  

\bibitem{bib-bib20} 
\par\noindent    Gaëtan Gilbert, Jesper Cockx, Matthieu Sozeau, and Nicolas Tabareau.  Definitional proof-irrelevance without K.  Proc. ACM Program. Lang., 3(POPL):3:1–3:28, 2019.  doi:10.1145/3290316.  

\bibitem{bib-bib21} 
\par\noindent    Jason Gross, Adam Chlipala, and David I. Spivak.  Experience Implementing a Performant Category-Theory Library in Coq.  In Gerwin Klein and Ruben Gamboa, editors, Interactive Theorem Proving - 5th International Conference, ITP 2014, Held as Part of the Vienna Summer of Logic, VSL 2014, Vienna, Austria, July 14-17, 2014. Proceedings, volume 8558 of Lecture Notes in Computer Science, pages 275–291. Springer, 2014.  doi:10.1007/978-3-319-08970-6\_18.  

\bibitem{bib-bib22} 
\par\noindent    Carl A. Gunter.  Universal Profinite Domains.  Inf. Comput., 72(1):1–30, 1987.  doi:10.1016/0890-5401(87)90048-4.  

\bibitem{bib-bib23} 
\par\noindent    Michael Hedberg.  A coherence theorem for martin-löf’s type theory.  J. Funct. Program., 8(4):413–436, 1998.  doi:10.1017/S0956796898003153.  

\bibitem{bib-bib24} 
\par\noindent    Jason Z. S. Hu and Jacques Carette.  Formalizing category theory in Agda.  In Catalin Hritcu and Andrei Popescu, editors, CPP ’21: 10th ACM SIGPLAN International Conference on Certified Programs and Proofs, Virtual Event, Denmark, January 17-19, 2021, pages 327–342. ACM, 2021.  doi:10.1145/3437992.3439922.  

\bibitem{bib-bib25} 
\par\noindent    Gérard P. Huet and Amokrane Saïbi.  Constructive category theory.  In Gordon D. Plotkin, Colin Stirling, and Mads Tofte, editors, Proof, Language, and Interaction, Essays in Honour of Robin Milner, pages 239–276. The MIT Press, 2000.  

\bibitem{bib-bib26} 
\par\noindent    Brian Huffman.  A Purely Definitional Universal Domain.  In Stefan Berghofer, Tobias Nipkow, Christian Urban, and Makarius Wenzel, editors, Theorem Proving in Higher Order Logics, 22nd International Conference, TPHOLs 2009, Munich, Germany, August 17-20, 2009. Proceedings, volume 5674 of Lecture Notes in Computer Science, pages 260–275. Springer, 2009.  doi:10.1007/978-3-642-03359-9\_19.  

\bibitem{bib-bib27} 
\par\noindent    Thomas Jech.  Set Theory.  Springer Berlin Heidelberg, 2003.  doi:10.1007/3-540-44761-x.  

\bibitem{bib-bib28} 
\par\noindent    Ralf Jung, Robbert Krebbers, Lars Birkedal, and Derek Dreyer.  Higher-order ghost state.  In Jacques Garrigue, Gabriele Keller, and Eijiro Sumii, editors, Proceedings of the 21st ACM SIGPLAN International Conference on Functional Programming, ICFP 2016, Nara, Japan, September 18-22, 2016, pages 256–269. ACM, 2016.  doi:10.1145/2951913.2951943.  

\bibitem{bib-bib29} 
\par\noindent    Ralf Jung, Robbert Krebbers, Jacques-Henri Jourdan, Ales Bizjak, Lars Birkedal, and Derek Dreyer.  Iris from the ground up: A modular foundation for higher-order concurrent separation logic.  J. Funct. Program., 28:e20, 2018.  doi:10.1017/S0956796818000151.  

\bibitem{bib-bib30} 
\par\noindent    Max Kelly.  Basic concepts of enriched category theory series number 64.  London Mathematical Society lecture note series. Cambridge University Press, Cambridge, England, February 1982.  

\bibitem{bib-bib31} 
\par\noindent    Dominik Kirst and Gert Smolka.  Large model constructions for second-order ZF in dependent type theory.  In June Andronick and Amy P. Felty, editors, Proceedings of the 7th ACM SIGPLAN International Conference on Certified Programs and Proofs, CPP 2018, Los Angeles, CA, USA, January 8-9, 2018, pages 228–239. ACM, 2018.  doi:10.1145/3167095.  

\bibitem{bib-bib32} 
\par\noindent    Robbert Krebbers, Jacques-Henri Jourdan, Ralf Jung, Joseph Tassarotti, Jan-Oliver Kaiser, Amin Timany, Arthur Charguéraud, and Derek Dreyer.  Mosel: a general, extensible modal framework for interactive proofs in separation logic.  Proc. ACM Program. Lang., 2(ICFP):77:1–77:30, 2018.  doi:10.1145/3236772.  

\bibitem{bib-bib33} 
\par\noindent    Robbert Krebbers, Amin Timany, and Lars Birkedal.  Interactive proofs in higher-order concurrent separation logic.  In Giuseppe Castagna and Andrew D. Gordon, editors, Proceedings of the 44th ACM SIGPLAN Symposium on Principles of Programming Languages, POPL 2017, Paris, France, January 18-20, 2017, pages 205–217. ACM, 2017.  doi:10.1145/3009837.3009855.  

\bibitem{bib-bib34} 
\par\noindent    Saunders Mac Lane.  Categories for the working mathematician.  Graduate Texts in Mathematics. Springer, New York, NY, 2 edition, September 1998.  

\bibitem{bib-bib35} 
\par\noindent    Saunders MacLane and Ieke Moerdijk.  Sheaves in geometry and logic.  Universitext. Springer, New York, NY, 1 edition, May 1992.  

\bibitem{bib-bib36} 
\par\noindent    The mathlib Community.  The lean mathematical library.  In Proceedings of the 9th ACM SIGPLAN International Conference on Certified Programs and Proofs, CPP 2020, pages 367–381, New York, NY, USA, 2020. Association for Computing Machinery.  doi:10.1145/3372885.3373824.  

\bibitem{bib-bib37} 
\par\noindent    Hiroshi Nakano.  A Modality for Recursion.  In 15th Annual IEEE Symposium on Logic in Computer Science, Santa Barbara, California, USA, June 26-29, 2000, pages 255–266. IEEE Computer Society, 2000.  doi:10.1109/LICS.2000.855774.  

\bibitem{bib-bib38} 
\par\noindent    Christine Paulin-Mohring.  A constructive denotational semantics for Kahn networks in Coq.  In Yves Bertot, Gérard Huet, Jean-Jacques Lévy, and Gordon Plotkin, editors, From Semantics to Computer Science, pages 383–413. Cambridge University Press, 2009.  URL: https://inria.hal.science/inria-00431806.  

\bibitem{bib-bib39} 
\par\noindent    Marco Paviotti, Rasmus Ejlers Møgelberg, and Lars Birkedal.  A model of PCF in guarded type theory.  In Dan R. Ghica, editor, The 31st Conference on the Mathematical Foundations of Programming Semantics, MFPS 2015, Nijmegen, The Netherlands, June 22-25, 2015, volume 319 of Electronic Notes in Theoretical Computer Science, pages 333–349. Elsevier, 2015.  doi:10.1016/J.ENTCS.2015.12.020.  

\bibitem{bib-bib40} 
\par\noindent    Daniel Peebles, James Deikun, Ulf Norell, Dan Doel, Darius Jahandarie, and James Cook.  categories: Categories parametrized by morphism equality in Agda.  https://github.com/copumpkin/categories, 2018.  

\bibitem{bib-bib41} 
\par\noindent    Dana Scott.  Outline of a mathematical theory of computation.  Technical Report PRG02, OUCL, November 1970.  

\bibitem{bib-bib42} 
\par\noindent    Dana Scott.  Continuous lattices.  In F. W. Lawvere, editor, Toposes, Algebraic Geometry and Logic, pages 97–136, Berlin, Heidelberg, 1972. Springer Berlin Heidelberg.  

\bibitem{bib-bib43} 
\par\noindent    Filip Sieczkowski, Ales Bizjak, and Lars Birkedal.  ModuRes: A Coq Library for Modular Reasoning About Concurrent Higher-Order Imperative Programming Languages.  In Christian Urban and Xingyuan Zhang, editors, Interactive Theorem Proving - 6th International Conference, ITP 2015, Nanjing, China, August 24-27, 2015, Proceedings, volume 9236 of Lecture Notes in Computer Science, pages 375–390. Springer, 2015.  doi:10.1007/978-3-319-22102-1\_25.  

\bibitem{bib-bib44} 
\par\noindent    Michael B. Smyth and Gordon D. Plotkin.  The Category-Theoretic Solution of Recursive Domain Equations.  SIAM J. Comput., 11(4):761–783, 1982.  doi:10.1137/0211062.  

\bibitem{bib-bib45} 
\par\noindent    Simon Spies, Lennard Gäher, Daniel Gratzer, Joseph Tassarotti, Robbert Krebbers, Derek Dreyer, and Lars Birkedal.  Transfinite Iris: resolving an existential dilemma of step-indexed separation logic.  In Stephen N. Freund and Eran Yahav, editors, PLDI ’21: 42nd ACM SIGPLAN International Conference on Programming Language Design and Implementation, Virtual Event, Canada, June 20-25, 2021, pages 80–95. ACM, 2021.  doi:10.1145/3453483.3454031.  

\bibitem{bib-bib46} 
\par\noindent    Sergei Stepanenko and Amin Timany.  The Rocq Mechanization of Solving Guarded Domain Equations in Presheaves Over Ordinals.  https://doi.org/10.5281/zenodo.15406039.  

\bibitem{bib-bib47} 
\par\noindent    V. Stoltenberg-Hansen, I. Lindström, and E. R. Griffor.  Mathematical Theory of Domains.  Cambridge Tracts in Theoretical Computer Science. Cambridge University Press, 1994.  

\bibitem{bib-bib48} 
\par\noindent    Amin Timany.  Commuting Quantifiers, October 2020.  Blog post, accessed on Feb 8, 2025.  URL: https://cs.au.dk/˜timany/blog/commuting\_quantifiers/.  

\bibitem{bib-bib49} 
\par\noindent    Amin Timany and Bart Jacobs.  Category Theory in Coq 8.5.  In Delia Kesner and Brigitte Pientka, editors, 1st International Conference on Formal Structures for Computation and Deduction (FSCD 2016), volume 52 of Leibniz International Proceedings in Informatics (LIPIcs), pages 30:1–30:18, Dagstuhl, Germany, 2016. Schloss Dagstuhl – Leibniz-Zentrum für Informatik.  doi:10.4230/LIPIcs.FSCD.2016.30.  

\bibitem{bib-bib50} 
\par\noindent    The Univalent Foundations Program.  Homotopy Type Theory: Univalent Foundations of Mathematics.  https://homotopytypetheory.org/book, Institute for Advanced Study, 2013.  

\bibitem{bib-bib51} 
\par\noindent    Mitchell Wand.  Fixed-Point Constructions in Order-Enriched Categories.  Theor. Comput. Sci., 8:13–30, 1979.  doi:10.1016/0304-3975(79)90053-7.  

\bibitem{bib-bib52} 
\par\noindent    John Wiegley.  category-theory: Category Theory in Coq.  https://github.com/jwiegley/category-theory, 2019.  

\end{thebibliography}
